\documentclass[a4paper,12pt]{article}
\usepackage[utf8]{inputenc}
\usepackage[german]{babel}
\usepackage{geometry}
\usepackage{graphicx}
\usepackage{tabularx}
\usepackage{booktabs}
\usepackage{xcolor}
\usepackage{hyperref}
\usepackage{fancyhdr}
\usepackage{enumitem}
\usepackage{listings}
\usepackage{tikz}
\usepackage{pgfgantt}
\usepackage{amsmath}
\usepackage{amssymb}

\geometry{margin=2.5cm}
\setlength{\parindent}{0pt}
\setlength{\parskip}{6pt}

% Code-Listing Konfiguration
\lstset{
    backgroundcolor=\color{gray!10},
    basicstyle=\ttfamily\footnotesize,
    breakatwhitespace=false,
    breaklines=true,
    captionpos=b,
    commentstyle=\color{green!50!black},
    frame=single,
    keywordstyle=\color{blue},
    language=Java,
    numbers=left,
    numberstyle=\tiny\color{gray},
    showspaces=false,
    showstringspaces=false,
    stringstyle=\color{red},
    tabsize=2,
    xleftmargin=2em,
    framexleftmargin=2em
}

% Header and Footer
\pagestyle{fancy}
\fancyhf{}
\fancyhead[L]{CIC to Give - Pflichtenheft}
\fancyhead[R]{\today}
\fancyfoot[C]{\thepage}

% Title page setup
\title{\Huge\textbf{Lehrabschlussprüfung}\\[0.5cm]\LARGE Pflichtenheft\\[0.3cm]\large CIC to Give - Fullstack Webanwendung\\[0.5cm]\normalsize Kostenlose Tauschbörse mit Spring Boot \& Next.js}
\author{Kandidat: Jorgic Petar\\Lehrberuf: Applikationsentwicklung - Coding\}
\date{\today}

\begin{document}

\maketitle
\newpage

\tableofcontents
\newpage

\section{Einleitung und Projektübersicht}

\subsection{Was ist ein Lasten- vs. Pflichtenheft?}

\textbf{Lastenheft:} Beschreibt \textit{WAS} der Auftraggeber möchte. Es definiert die Anforderungen aus Sicht des Kunden.

\textbf{Pflichtenheft:} Beschreibt \textit{WIE} der Auftragnehmer (Entwickler) die Anforderungen technisch umsetzt.

Dieses Dokument kombiniert beide Aspekte für die Lehrabschlussprüfung.

\subsection{Projektbeschreibung}
CIC to Give ist eine vollständige Fullstack-Webanwendung, die als kostenlose Tauschbörse fungiert. Benutzer können Items einstellen, die sie verschenken möchten, und nach Items suchen, die sie benötigen.

\subsection{Warum diese Technologien?}

\textbf{Java mit Spring Boot:}
\begin{itemize}
    \item Typsicher und robust
    \item Große Enterprise-Unterstützung
    \item Viele vorgefertigte Lösungen (Spring Security, Spring Data)
    \item Gute Performance für Backend-Services
\end{itemize}

\textbf{Next.js mit React:}
\begin{itemize}
    \item Server-Side Rendering für bessere SEO
    \item Moderne Frontend-Architektur
    \item TypeScript-Unterstützung
    \item Große Community und Dokumentation
\end{itemize}

\section{Technische Architektur}

\subsection{System-Architektur}

\begin{center}
\begin{tikzpicture}[node distance=2cm]
    \node[rectangle, draw, fill=blue!20] (frontend) {Next.js Frontend\\(Port 3000)};
    \node[rectangle, draw, fill=green!20, below of=frontend] (backend) {Spring Boot Backend\\(Port 8080)};
    \node[rectangle, draw, fill=yellow!20, below of=backend] (database) {PostgreSQL Database\\(Port 5432)};

    \draw[->] (frontend) -- (backend) node[midway, right] {HTTP/REST API};
    \draw[->] (backend) -- (database) node[midway, right] {JDBC};
\end{tikzpicture}
\end{center}

\subsection{Warum diese Architektur?}

\textbf{Drei-Schichten-Architektur:}
\begin{enumerate}
    \item \textbf{Präsentationsschicht:} Next.js Frontend
    \item \textbf{Anwendungsschicht:} Spring Boot Backend
    \item \textbf{Datenschicht:} PostgreSQL Database
\end{enumerate}

\textbf{Vorteile:}
\begin{itemize}
    \item Klare Trennung der Verantwortlichkeiten
    \item Skalierbarkeit (jede Schicht kann separat skaliert werden)
    \item Wartbarkeit (Änderungen in einer Schicht betreffen andere nicht)
    \item Testbarkeit (jede Schicht kann isoliert getestet werden)
\end{itemize}

\section{Backend-Implementierung mit Spring Boot}

\subsection{Warum Spring Boot?}

Spring Boot ist ein Framework für Java, das die Entwicklung von Enterprise-Anwendungen vereinfacht:

\begin{itemize}
    \item \textbf{Auto-Configuration:} Automatische Konfiguration basierend auf Classpath
    \item \textbf{Embedded Server:} Tomcat ist bereits integriert
    \item \textbf{Starter Dependencies:} Vorgefertigte Dependency-Pakete
    \item \textbf{Production-Ready:} Metrics, Health Checks, etc. sind bereits implementiert
\end{itemize}

\subsection{Maven Konfiguration}

\begin{lstlisting}[language=XML, caption=pom.xml - Maven Dependencies]
<dependencies>
    <!-- Spring Boot Web Starter -->
    <dependency>
        <groupId>org.springframework.boot</groupId>
        <artifactId>spring-boot-starter-web</artifactId>
    </dependency>

    <!-- Spring Boot Data JPA -->
    <dependency>
        <groupId>org.springframework.boot</groupId>
        <artifactId>spring-boot-starter-data-jpa</artifactId>
    </dependency>

    <!-- Spring Security -->
    <dependency>
        <groupId>org.springframework.boot</groupId>
        <artifactId>spring-boot-starter-security</artifactId>
    </dependency>

    <!-- JWT Token Support -->
    <dependency>
        <groupId>io.jsonwebtoken</groupId>
        <artifactId>jjwt-api</artifactId>
        <version>0.11.5</version>
    </dependency>
</dependencies>
\end{lstlisting}

\textbf{Warum diese Dependencies?}
\begin{itemize}
    \item \textbf{spring-boot-starter-web:} Enthält Spring MVC, Tomcat, Jackson für JSON
    \item \textbf{spring-boot-starter-data-jpa:} Hibernate ORM, Spring Data JPA
    \item \textbf{spring-boot-starter-security:} Authentifizierung und Autorisierung
    \item \textbf{jjwt:} JWT Token-Generierung und -Validierung
\end{itemize}

\subsection{Datenmodell (JPA Entities)}

\subsubsection{User Entity}

\begin{lstlisting}[language=Java, caption=User.java - Benutzer-Entität]
@Entity
@Table(name = "users",
       uniqueConstraints = {
           @UniqueConstraint(columnNames = "username"),
           @UniqueConstraint(columnNames = "email")
       })
public class User {

    @Id
    @GeneratedValue(strategy = GenerationType.IDENTITY)
    private Long id;

    @Column(nullable = false, length = 50)
    private String username;

    @Column(nullable = false, length = 100)
    private String email;

    @Column(nullable = false)
    private String password; // BCrypt gehashed

    @OneToMany(mappedBy = "user", cascade = CascadeType.ALL, fetch = FetchType.LAZY)
    @JsonIgnore
    private Set<Item> items = new HashSet<>();

    // Konstruktoren, Getter, Setter...
}
\end{lstlisting}

\textbf{Erklärung der Annotationen:}
\begin{itemize}
    \item \textbf{@Entity:} Markiert die Klasse als JPA-Entität
    \item \textbf{@Table:} Definiert Tabellennamen und Constraints
    \item \textbf{@Id:} Primärschlüssel
    \item \textbf{@GeneratedValue:} Auto-Increment für ID
    \item \textbf{@Column:} Spaltendefinition (nullable, length)
    \item \textbf{@OneToMany:} Ein User kann viele Items haben
    \item \textbf{@UniqueConstraint:} Username und Email müssen einzigartig sein
\end{itemize}

\subsubsection{Item Entity}

\begin{lstlisting}[language=Java, caption=Item.java - Item-Entität]
@Entity
@JsonIgnoreProperties({"hibernateLazyInitializer", "handler"})
public class Item {

    @Id
    @GeneratedValue(strategy = GenerationType.IDENTITY)
    private Long id;

    @Column(nullable = false)
    private String name;

    @Column(length = 2000)
    private String description;

    @Column(name = "created_at")
    private LocalDateTime createdAt;

    @ElementCollection
    @CollectionTable(
        name = "item_images",
        joinColumns = @JoinColumn(name = "item_id")
    )
    @Column(name = "image_url")
    private List<String> imageUrls = new ArrayList<>();

    @ManyToOne(fetch = FetchType.LAZY)
    @JoinColumn(name = "user_id")
    @JsonIgnoreProperties({"password", "items"})
    private User user;

    @Column(nullable = false)
    private String category;

    @Column(nullable = false)
    private boolean reserved = false;

    @OneToMany(mappedBy = "item", cascade = CascadeType.ALL, orphanRemoval = true)
    @JsonIgnoreProperties({"item"})
    private List<Comment> comments = new ArrayList<>();

    @PreUpdate
    public void preUpdate() {
        this.updatedAt = LocalDateTime.now();
    }
}
\end{lstlisting}

\textbf{Komplexere Annotationen erklärt:}
\begin{itemize}
    \item \textbf{@ElementCollection:} Für die Liste von Bild-URLs
    \item \textbf{@CollectionTable:} Definiert separate Tabelle für Bilder
    \item \textbf{@ManyToOne:} Viele Items gehören zu einem User
    \item \textbf{@OneToMany:} Ein Item kann viele Kommentare haben
    \item \textbf{@PreUpdate:} Wird vor jedem Update ausgeführt
    \item \textbf{@JsonIgnoreProperties:} Verhindert JSON-Serialisierung bestimmter Felder
\end{itemize}

\subsection{Passwort-Hashing mit BCrypt}

\subsubsection{Warum BCrypt?}

\textbf{Problem:} Passwörter dürfen niemals im Klartext gespeichert werden!

\textbf{Lösung:} BCrypt ist ein sicherer Hashing-Algorithmus:
\begin{itemize}
    \item \textbf{Salzung:} Jedes Passwort bekommt einen zufälligen Salt
    \item \textbf{Adaptive Hashing:} Rechenaufwand kann erhöht werden
    \item \textbf{Einweg-Funktion:} Aus dem Hash kann nicht das Original-Passwort berechnet werden
\end{itemize}

\subsubsection{Technische Implementierung}

\begin{lstlisting}[language=Java, caption=Passwort-Hashing in Spring Security]
@Configuration
@EnableWebSecurity
public class WebSecurityConfig {

    @Bean
    public PasswordEncoder passwordEncoder() {
        return new BCryptPasswordEncoder();
    }

    // Passwort beim Registrieren hashen
    public String hashPassword(String plainPassword) {
        return passwordEncoder().encode(plainPassword);
    }

    // Passwort beim Login prüfen
    public boolean checkPassword(String plainPassword, String hashedPassword) {
        return passwordEncoder().matches(plainPassword, hashedPassword);
    }
}
\end{lstlisting}

\textbf{Wie funktioniert BCrypt?}
\begin{enumerate}
    \item \textbf{Salt generieren:} Zufällige 128-bit Zeichenkette
    \item \textbf{Passwort + Salt:} Kombiniert das Passwort mit dem Salt
    \item \textbf{Hashing:} Verwendet Blowfish-Algorithmus mit konfigurierbaren Runden
    \item \textbf{Speichern:} Salt und Hash werden zusammen gespeichert
\end{enumerate}

\textbf{Beispiel:}
\begin{itemize}
    \item \textbf{Original:} "meinPasswort123"
    \item \textbf{BCrypt Hash:}
    "\$2a\$10\$N9qo8uLOickgx2ZMRZoMye6FaQJdFfhWO/FdJwF7dWQbN8A5JzNQK"
    \item \textbf{Enthält:} Algorithmus-Info, Kosten-Parameter, Salt und Hash
\end{itemize}

\subsection{JWT (JSON Web Token) Authentifizierung}

\subsubsection{Warum JWT statt Sessions?}

\textbf{Traditionelle Sessions:}
\begin{itemize}
    \item Server speichert Session-Daten im Speicher
    \item Session-ID wird als Cookie übertragen
    \item Problem: Nicht skalierbar bei mehreren Servern
\end{itemize}

\textbf{JWT Vorteile:}
\begin{itemize}
    \item \textbf{Stateless:} Server muss keine Session-Daten speichern
    \item \textbf{Skalierbar:} Funktioniert mit Load Balancern
    \item \textbf{Cross-Domain:} Kann zwischen verschiedenen Domains verwendet werden
    \item \textbf{Selbst-validierend:} Token enthält alle nötigen Informationen
\end{itemize}

\subsubsection{JWT Struktur}

Ein JWT besteht aus drei Teilen, getrennt durch Punkte:
\textbf{Header.Payload.Signature}

\begin{lstlisting}[language=JSON, caption=JWT Header (Base64 encoded)]
{
  "alg": "HS256",
  "typ": "JWT"
}
\end{lstlisting}

\begin{lstlisting}[language=JSON, caption=JWT Payload (Base64 encoded)]
{
  "sub": "1234567890",
  "username": "john_doe",
  "iat": 1516239022,
  "exp": 1516325422
}
\end{lstlisting}

\textbf{Signature:} HMACSHA256(base64UrlEncode(header) + "." + base64UrlEncode(payload), secret)

\subsubsection{JWT Implementierung}

\begin{lstlisting}[language=Java, caption=JwtUtils.java - JWT Token-Generierung]
@Component
public class JwtUtils {

    private String jwtSecret = "mySecretKey";
    private int jwtExpirationMs = 86400000; // 24 Stunden

    public String generateJwtToken(UserDetailsImpl userPrincipal) {
        return generateTokenFromUsername(userPrincipal.getUsername());
    }

    public String generateTokenFromUsername(String username) {
        return Jwts.builder()
                .setSubject(username)
                .setIssuedAt(new Date())
                .setExpiration(new Date((new Date()).getTime() + jwtExpirationMs))
                .signWith(SignatureAlgorithm.HS512, jwtSecret)
                .compact();
    }

    public String getUserNameFromJwtToken(String token) {
        return Jwts.parser()
                .setSigningKey(jwtSecret)
                .parseClaimsJws(token)
                .getBody()
                .getSubject();
    }

    public boolean validateJwtToken(String authToken) {
        try {
            Jwts.parser().setSigningKey(jwtSecret).parseClaimsJws(authToken);
            return true;
        } catch (SignatureException e) {
            logger.error("Invalid JWT signature: {}", e.getMessage());
        } catch (MalformedJwtException e) {
            logger.error("Invalid JWT token: {}", e.getMessage());
        } catch (ExpiredJwtException e) {
            logger.error("JWT token is expired: {}", e.getMessage());
        }
        return false;
    }
}
\end{lstlisting}

\subsection{Spring Security Konfiguration}

\subsubsection{Warum Spring Security?}

Spring Security ist das De-facto-Standard-Framework für Sicherheit in Java-Anwendungen:
\begin{itemize}
    \item \textbf{Authentifizierung:} Wer ist der Benutzer?
    \item \textbf{Autorisierung:} Was darf der Benutzer?
    \item \textbf{Schutz vor Angriffen:} CSRF, XSS, Session Fixation
\end{itemize}

\begin{lstlisting}[language=Java, caption=WebSecurityConfig.java - Security-Konfiguration]
@Configuration
@EnableWebSecurity
public class WebSecurityConfig {

    @Autowired
    private JwtAuthenticationEntryPoint unauthorizedHandler;

    @Bean
    public JwtAuthenticationFilter authenticationJwtTokenFilter() {
        return new JwtAuthenticationFilter();
    }

    @Bean
    public PasswordEncoder passwordEncoder() {
        return new BCryptPasswordEncoder();
    }

    @Bean
    public SecurityFilterChain filterChain(HttpSecurity http) throws Exception {
        http.cors().and().csrf().disable()
            .exceptionHandling().authenticationEntryPoint(unauthorizedHandler).and()
            .sessionManagement().sessionCreationPolicy(SessionCreationPolicy.STATELESS).and()
            .authorizeHttpRequests()
                .requestMatchers("/api/auth/**").permitAll()
                .requestMatchers("/api/items/**").authenticated()
                .anyRequest().authenticated();

        http.addFilterBefore(authenticationJwtTokenFilter(), UsernamePasswordAuthenticationFilter.class);

        return http.build();
    }
}
\end{lstlisting}

\textbf{Konfiguration erklärt:}
\begin{itemize}
    \item \textbf{cors().and().csrf().disable():} CORS aktivieren, CSRF deaktivieren (für APIs)
    \item \textbf{sessionManagement().sessionCreationPolicy(STATELESS):} Keine Sessions verwenden
    \item \textbf{authorizeHttpRequests():} Definiert welche URLs geschützt sind
    \item \textbf{addFilterBefore():} JWT-Filter vor Standard-Filter einfügen
\end{itemize}

\subsection{REST API Design}

\subsubsection{Was ist REST?}

REST (Representational State Transfer) ist ein Architekturstil für APIs:
\begin{itemize}
    \item \textbf{Stateless:} Jede Anfrage enthält alle nötigen Informationen
    \item \textbf{HTTP-Verben:} GET, POST, PUT, DELETE für verschiedene Operationen
    \item \textbf{Ressourcen:} URLs repräsentieren Objekte
    \item \textbf{JSON:} Standardformat für Datenübertragung
\end{itemize}

\subsubsection{ItemController Implementierung}

\begin{lstlisting}[language=Java, caption=ItemController.java - REST-Controller]
@RestController
@RequestMapping("/api/items")
@CrossOrigin(origins = "http://localhost:3000")
public class ItemController {

    @Autowired
    private ItemRepository itemRepository;

    @Autowired
    private UserRepository userRepository;

    @GetMapping
    public ResponseEntity<Map<String, Object>> getAllItems(
            @RequestParam(defaultValue = "0") int page,
            @RequestParam(defaultValue = "10") int size,
            @RequestParam(defaultValue = "createdAt") String sortBy,
            @RequestParam(defaultValue = "desc") String sortDir,
            @RequestParam(required = false) String category,
            @RequestParam(required = false) String search) {

        Sort sort = sortDir.equalsIgnoreCase("desc") ?
            Sort.by(sortBy).descending() :
            Sort.by(sortBy).ascending();

        Pageable pageable = PageRequest.of(page, size, sort);
        Page<Item> itemPage = itemRepository.findItemsWithFilters(category, search, pageable);

        Map<String, Object> response = new HashMap<>();
        response.put("items", itemPage.getContent());
        response.put("currentPage", itemPage.getNumber());
        response.put("totalItems", itemPage.getTotalElements());
        response.put("totalPages", itemPage.getTotalPages());

        return ResponseEntity.ok(response);
    }

    @PostMapping
    public ResponseEntity<Item> createItem(@RequestBody Item item, Authentication authentication) {
        String username = authentication.getName();
        User user = userRepository.findByUsername(username)
            .orElseThrow(() -> new RuntimeException("User not found"));

        item.setUser(user);
        Item savedItem = itemRepository.save(item);
        return ResponseEntity.status(HttpStatus.CREATED).body(savedItem);
    }
}
\end{lstlisting}

\textbf{Annotationen erklärt:}
\begin{itemize}
    \item \textbf{@RestController:} Kombiniert @Controller und @ResponseBody
    \item \textbf{@RequestMapping:} Basis-URL für alle Methoden
    \item \textbf{@GetMapping:} HTTP GET-Requests
    \item \textbf{@PostMapping:} HTTP POST-Requests
    \item \textbf{@RequestParam:} URL-Parameter
    \item \textbf{@RequestBody:} JSON-Body in Java-Objekt konvertieren
\end{itemize}

\subsection{Datenbankzugriff mit Spring Data JPA}

\subsubsection{Warum JPA?}

JPA (Java Persistence API) ist der Standard für Object-Relational Mapping in Java:
\begin{itemize}
    \item \textbf{Objektrelational:} Java-Objekte werden auf Datenbanktabellen gemappt
    \item \textbf{Provider-unabhängig:} Funktioniert mit verschiedenen Datenbanken
    \item \textbf{Hibernate:} Meist verwendete JPA-Implementierung
\end{itemize}

\begin{lstlisting}[language=Java, caption=ItemRepository.java - Datenbankzugriff]
@Repository
public interface ItemRepository extends JpaRepository<Item, Long> {

    @Query("SELECT i FROM Item i WHERE " +
           "(:category IS NULL OR i.category = :category) AND " +
           "(:search IS NULL OR LOWER(i.name) LIKE LOWER(CONCAT('%', :search, '%')) OR " +
           "LOWER(i.description) LIKE LOWER(CONCAT('%', :search, '%')))")
    Page<Item> findItemsWithFilters(@Param("category") String category,
                                   @Param("search") String search,
                                   Pageable pageable);

    List<Item> findByUserOrderByCreatedAtDesc(User user);

    @Query("SELECT i FROM Item i WHERE i.user.id = :userId")
    List<Item> findByUserId(@Param("userId") Long userId);
}
\end{lstlisting}

\textbf{Spring Data JPA Features:}
\begin{itemize}
    \item \textbf{Automatische Implementierung:} Spring erstellt die Implementierung zur Laufzeit
    \item \textbf{Query Methods:} Methodennamen werden zu SQL-Queries
    \item \textbf{@Query:} Eigene JPQL-Queries definieren
    \item \textbf{Paginierung:} Pageable für große Datenmengen
\end{itemize}

\section{Frontend-Implementierung mit Next.js}

\subsection{Warum Next.js?}

Next.js ist ein React-Framework mit zusätzlichen Features:
\begin{itemize}
    \item \textbf{Server-Side Rendering:} Bessere Performance und SEO
    \item \textbf{File-based Routing:} Ordnerstruktur definiert URLs
    \item \textbf{API Routes:} Backend-Funktionalität im Frontend-Projekt
    \item \textbf{Automatic Code Splitting:} Optimierte Bundle-Größen
\end{itemize}

\subsection{TypeScript Integration}

\subsubsection{Warum TypeScript?}

TypeScript ist JavaScript mit Typen:
\begin{itemize}
    \item \textbf{Typsicherheit:} Fehler werden zur Compile-Zeit erkannt
    \item \textbf{Bessere IDE-Unterstützung:} Autocompletion, Refactoring
    \item \textbf{Selbstdokumentierend:} Typen sind Dokumentation
    \item \textbf{Großprojekt-tauglich:} Bessere Wartbarkeit bei vielen Entwicklern
\end{itemize}

\begin{lstlisting}[language=TypeScript, caption=api.ts - TypeScript-Interfaces]
// Type definitions
export interface User {
  id: number;
  username: string;
  email: string;
}

export interface Item {
  id?: number;
  name: string;
  description: string;
  category: string;
  location?: string;
  imageUrls?: string[];
  createdAt?: string;
  updatedAt?: string;
  user?: User;
  reserved?: boolean;
}

export interface LoginRequest {
  username: string;
  password: string;
}

export interface PaginatedResponse<T> {
  items: T[];
  currentPage: number;
  totalItems: number;
  totalPages: number;
}
\end{lstlisting}

\subsection{API-Service Implementierung}

\begin{lstlisting}[language=TypeScript, caption=API-Service für Backend-Kommunikation]
const API_BASE_URL = "http://localhost:8080/api";

// Helper function to get auth header
const getAuthHeader = (): HeadersInit => {
  const token = localStorage.getItem("authToken");
  return token ? { Authorization: `Bearer ${token}` } : {};
};

export const authApi = {
  login: async (loginRequest: LoginRequest): Promise<LoginResponse> => {
    const response = await fetch(`${API_BASE_URL}/auth/login`, {
      method: "POST",
      headers: {
        "Content-Type": "application/json",
        Accept: "application/json",
      },
      credentials: "include",
      body: JSON.stringify(loginRequest),
    });

    if (!response.ok) {
      throw new Error(`HTTP ${response.status}: Login failed`);
    }

    const data = await response.json();

    // Store user info und token
    localStorage.setItem("user", JSON.stringify(data));
    localStorage.setItem("isLoggedIn", "true");

    if (data.token) {
      localStorage.setItem("authToken", data.token);
    }

    return data;
  },

  isLoggedIn: (): boolean => {
    if (typeof window !== "undefined") {
      const isLoggedIn = localStorage.getItem("isLoggedIn") === "true";
      const hasUser = !!localStorage.getItem("user");
      return isLoggedIn && hasUser;
    }
    return false;
  }
};
\end{lstlisting}

\textbf{Frontend-Backend-Kommunikation:}
\begin{itemize}
    \item \textbf{Fetch API:} Moderne Alternative zu XMLHttpRequest
    \item \textbf{localStorage:} Client-seitige Speicherung von Tokens
    \item \textbf{Authorization Header:} JWT-Token in jedem Request
    \item \textbf{Error Handling:} Strukturierte Fehlerbehandlung
\end{itemize}

\subsection{React Komponenten-Architektur}

\begin{lstlisting}[language=TypeScript, caption=ItemList.tsx - React-Komponente]
"use client";
import { useEffect, useState } from "react";
import { itemsApi, Item } from "@/services/api";

interface ItemListProps {
  category?: string;
  search?: string;
}

export default function ItemList({ category, search }: ItemListProps) {
  const [items, setItems] = useState<Item[]>([]);
  const [loading, setLoading] = useState(true);
  const [error, setError] = useState<string | null>(null);

  useEffect(() => {
    const fetchItems = async () => {
      setLoading(true);
      setError(null);

      try {
        const response = await itemsApi.getAll(0, 10, category, search);
        setItems(response.items);
      } catch (err) {
        setError(err instanceof Error ? err.message : "Fehler beim Laden");
      } finally {
        setLoading(false);
      }
    };

    fetchItems();
  }, [category, search]);

  if (loading) return <div>Laden...</div>;
  if (error) return <div>Fehler: {error}</div>;

  return (
    <div className="grid grid-cols-1 md:grid-cols-2 lg:grid-cols-3 gap-4">
      {items.map((item) => (
        <div key={item.id} className="border rounded-lg p-4">
          <h3 className="font-bold">{item.name}</h3>
          <p className="text-gray-600">{item.description}</p>
          <div className="mt-2">
            <span className="bg-blue-100 text-blue-800 px-2 py-1 rounded">
              {item.category}
            </span>
          </div>
        </div>
      ))}
    </div>
  );
}
\end{lstlisting}

\textbf{React Konzepte erklärt:}
\begin{itemize}
    \item \textbf{useState:} Hook für lokalen Komponenten-State
    \item \textbf{useEffect:} Hook für Side Effects (API-Calls, Event Listener)
    \item \textbf{Props:} Parameter, die zwischen Komponenten übergeben werden
    \item \textbf{JSX:} JavaScript mit HTML-ähnlicher Syntax
    \item \textbf{Component Lifecycle:} Mount, Update, Unmount
\end{itemize}

\subsection{Styling mit Tailwind CSS}

\subsubsection{Warum Tailwind CSS?}

Tailwind ist ein Utility-First CSS-Framework:
\begin{itemize}
    \item \textbf{Utility-Classes:} Vordefinierte CSS-Klassen für alle Properties
    \item \textbf{Responsive Design:} Breakpoints sind integriert
    \item \textbf{Konsistenz:} Design-System ist vorgegeben
    \item \textbf{Performance:} Nur verwendete CSS-Klassen werden gebündelt
\end{itemize}

\begin{lstlisting}[language=HTML, caption=Tailwind CSS Beispiel]
<div className="grid grid-cols-1 md:grid-cols-2 lg:grid-cols-3 gap-4 p-6">
  <div className="bg-white rounded-lg shadow-md p-4 hover:shadow-lg transition-shadow">
    <h3 className="text-lg font-bold text-gray-900 mb-2">Titel</h3>
    <p className="text-gray-600 text-sm line-clamp-3">Beschreibung...</p>
    <div className="mt-4 flex items-center justify-between">
      <span className="bg-green-100 text-green-800 px-2 py-1 rounded-full text-xs">
        Kategorie
      </span>
      <button className="bg-blue-500 hover:bg-blue-600 text-white px-4 py-2 rounded">
        Details
      </button>
    </div>
  </div>
</div>
\end{lstlisting}

\textbf{Tailwind-Klassen erklärt:}
\begin{itemize}
    \item \textbf{grid grid-cols-1:} CSS Grid mit 1 Spalte
    \item \textbf{md:grid-cols-2:} 2 Spalten ab Medium-Breakpoint
    \item \textbf{bg-white:} Weißer Hintergrund
    \item \textbf{rounded-lg:} Abgerundete Ecken
    \item \textbf{shadow-md:} Mittelstarker Schatten
    \item \textbf{hover:shadow-lg:} Stärkerer Schatten beim Hover
    \item \textbf{transition-shadow:} Animierte Schatten-Übergänge
\end{itemize}

\section{Datenbank-Design}

\subsection{Warum PostgreSQL?}

PostgreSQL ist eine objektrelationale Datenbank:
\begin{itemize}
    \item \textbf{ACID-Eigenschaften:} Atomarität, Konsistenz, Isolation, Dauerhaftigkeit
    \item \textbf{Skalierbarkeit:} Horizontal und vertikal skalierbar
    \item \textbf{JSON-Unterstützung:} Native JSON/JSONB Datentypen
    \item \textbf{Open Source:} Kostenlos und community-supported
    \item \textbf{Erweiterbarkeit:} Custom Functions, Trigger, Extensions
\end{itemize}

\subsection{Datenbankschema}

\begin{center}
\begin{tikzpicture}[scale=0.8]
    % Users Table
    \node[rectangle, draw, fill=blue!20] (users) at (0, 0) {
        \begin{tabular}{l}
        \textbf{users} \\
        \hline
        id (PK) \\
        username \\
        email \\
        password \\
        created\_at
        \end{tabular}
    };

    % Items Table
    \node[rectangle, draw, fill=green!20] (items) at (6, 0) {
        \begin{tabular}{l}
        \textbf{items} \\
        \hline
        id (PK) \\
        name \\
        description \\
        category \\
        location \\
        reserved \\
        user\_id (FK) \\
        created\_at \\
        updated\_at
        \end{tabular}
    };

    % Comments Table
    \node[rectangle, draw, fill=yellow!20] (comments) at (3, -4) {
        \begin{tabular}{l}
        \textbf{comments} \\
        \hline
        id (PK) \\
        text \\
        user\_id (FK) \\
        item\_id (FK) \\
        created\_at
        \end{tabular}
    };

    % Item Images Table
    \node[rectangle, draw, fill=red!20] (images) at (6, -2.5) {
        \begin{tabular}{l}
        \textbf{item\_images} \\
        \hline
        item\_id (FK) \\
        image\_url
        \end{tabular}
    };

    % Relationships
    \draw[->] (users) -- (items) node[midway, above] {1:n};
    \draw[->] (users) -- (comments) node[midway, left] {1:n};
    \draw[->] (items) -- (comments) node[midway, right] {1:n};
    \draw[->] (items) -- (images) node[midway, right] {1:n};
\end{tikzpicture}
\end{center}

\subsection{Datenbank-Konfiguration}

\begin{lstlisting}[language=Properties, caption=application.properties - Datenbankkonfiguration]
# H2 Database (Development)
spring.h2.console.enabled=true
spring.datasource.url=jdbc:h2:mem:testdb
spring.datasource.username=sa
spring.datasource.password=password

# PostgreSQL (Production)
spring.datasource.url=jdbc:postgresql://localhost:5432/cictogive
spring.datasource.username=${DB_USERNAME:postgres}
spring.datasource.password=${DB_PASSWORD:password}

# JPA Configuration
spring.jpa.hibernate.ddl-auto=update
spring.jpa.show-sql=true
spring.jpa.properties.hibernate.dialect=org.hibernate.dialect.PostgreSQLDialect

# File Upload Configuration
spring.servlet.multipart.max-file-size=10MB
spring.servlet.multipart.max-request-size=50MB
\end{lstlisting}

\textbf{Konfiguration erklärt:}
\begin{itemize}
    \item \textbf{ddl-auto=update:} Hibernate erstellt/aktualisiert Tabellen automatisch
    \item \textbf{show-sql=true:} SQL-Queries werden in der Konsole angezeigt
    \item \textbf{Environment Variables:} \${DB\_USERNAME:postgres} mit Fallback
    \item \textbf{Multipart:} Konfiguration für Datei-Upload
\end{itemize}

\section{Sicherheitskonzepte}

\subsection{Authentication vs. Authorization}

\textbf{Authentication (Authentifizierung):} \textit{Wer ist der Benutzer?}
\begin{itemize}
    \item Login mit Username/Passwort
    \item JWT-Token-Validierung
    \item Session-Management
\end{itemize}

\textbf{Authorization (Autorisierung):} \textit{Was darf der Benutzer?}
\begin{itemize}
    \item Benutzer darf nur eigene Items bearbeiten
    \item Bestimmte Endpoints sind nur für eingeloggte Benutzer
    \item Admin-Funktionen (falls vorhanden)
\end{itemize}

\subsection{CORS (Cross-Origin Resource Sharing)}

\textbf{Problem:} Browser blockieren Requests zwischen verschiedenen Domains aus Sicherheitsgründen.

\textbf{Lösung:} CORS-Header erlauben bestimmte Cross-Origin-Requests.

\begin{lstlisting}[language=Java, caption=CORS-Konfiguration]
@CrossOrigin(origins = "http://localhost:3000")
@RestController
public class ItemController {
    // Controller-Methoden...
}

// Oder global:
@Configuration
public class WebConfig implements WebMvcConfigurer {
    @Override
    public void addCorsMappings(CorsRegistry registry) {
        registry.addMapping("/api/**")
                .allowedOrigins("http://localhost:3000")
                .allowedMethods("GET", "POST", "PUT", "DELETE")
                .allowCredentials(true);
    }
}
\end{lstlisting}

\subsection{Input-Validierung}

\textbf{Warum wichtig?} Schutz vor SQL-Injection, XSS, und Dateninkonsistenz.

\begin{lstlisting}[language=Java, caption=Bean Validation]
public class Item {
    @NotBlank(message = "Name ist erforderlich")
    @Size(max = 100, message = "Name darf maximal 100 Zeichen haben")
    private String name;

    @Size(max = 2000, message = "Beschreibung darf maximal 2000 Zeichen haben")
    private String description;

    @NotBlank(message = "Kategorie ist erforderlich")
    private String category;

    @Email(message = "Ungültige E-Mail-Adresse")
    private String email;
}

@PostMapping
public ResponseEntity<Item> createItem(@Valid @RequestBody Item item) {
    // @Valid triggert automatische Validierung
    // Bei Fehlern wird automatisch 400 Bad Request zurückgegeben
}
\end{lstlisting}

\section{File Upload und Storage}

\subsection{Warum Datei-Upload komplex ist}

\begin{itemize}
    \item \textbf{Sicherheit:} Malware-Uploads verhindern
    \item \textbf{Performance:} Große Dateien blockieren nicht den Server
    \item \textbf{Storage:} Wo werden Dateien gespeichert?
    \item \textbf{Skalierung:} Was passiert bei Millionen von Bildern?
\end{itemize}

\subsection{Implementierung}

\begin{lstlisting}[language=Java, caption=FileUploadController.java]
@RestController
@RequestMapping("/api/upload")
public class FileUploadController {

    private final String uploadDir = "uploads/";

    @PostMapping("/images")
    public ResponseEntity<List<String>> uploadImages(
            @RequestParam("files") MultipartFile[] files) {

        List<String> imageUrls = new ArrayList<>();

        for (MultipartFile file : files) {
            // Validierung
            if (!isValidImageFile(file)) {
                throw new RuntimeException("Ungültiger Dateityp: " + file.getContentType());
            }

            if (file.getSize() > 10 * 1024 * 1024) { // 10MB
                throw new RuntimeException("Datei zu groß: " + file.getSize());
            }

            try {
                // Eindeutigen Dateinamen generieren
                String fileName = UUID.randomUUID().toString() + "_" + file.getOriginalFilename();
                Path targetPath = Paths.get(uploadDir + fileName);

                // Ordner erstellen falls nicht vorhanden
                Files.createDirectories(targetPath.getParent());

                // Datei speichern
                file.transferTo(targetPath);

                // URL für Frontend
                String imageUrl = "/uploads/" + fileName;
                imageUrls.add(imageUrl);

            } catch (IOException e) {
                throw new RuntimeException("Fehler beim Speichern: " + e.getMessage());
            }
        }

        return ResponseEntity.ok(imageUrls);
    }

    private boolean isValidImageFile(MultipartFile file) {
        String contentType = file.getContentType();
        return contentType != null &&
               (contentType.equals("image/jpeg") ||
                contentType.equals("image/png") ||
                contentType.equals("image/gif"));
    }
}
\end{lstlisting}

\textbf{Sicherheitsmaßnahmen:}
\begin{itemize}
    \item Dateityp-Validierung (nur Bilder)
    \item Dateigrößen-Limit (10MB)
    \item Eindeutige Dateinamen (UUID)
    \item Path Traversal-Schutz
\end{itemize}

\section{Testing}

\subsection{Warum Testen wichtig ist}

\begin{itemize}
    \item \textbf{Qualitätssicherung:} Bugs früh erkennen
    \item \textbf{Refactoring-Sicherheit:} Änderungen brechen nichts
    \item \textbf{Dokumentation:} Tests zeigen wie Code verwendet wird
    \item \textbf{Confidence:} Sicherheit bei Deployments
\end{itemize}

\subsection{Test-Pyramid}

\begin{center}
\begin{tikzpicture}
    \filldraw[fill=green!30] (0,0) -- (6,0) -- (5,1) -- (1,1) -- cycle;
    \filldraw[fill=yellow!30] (1,1) -- (5,1) -- (4,2) -- (2,2) -- cycle;
    \filldraw[fill=red!30] (2,2) -- (4,2) -- (3.5,2.5) -- (2.5,2.5) -- cycle;

    \node at (3,0.5) {\textbf{Unit Tests (viele)}};
    \node at (3,1.5) {\textbf{Integration Tests (weniger)}};
    \node at (3,2.25) {\textbf{E2E Tests (wenige)}};
\end{tikzpicture}
\end{center}

\subsection{Unit Tests}

\begin{lstlisting}[language=Java, caption=ItemControllerTest.java]
@SpringBootTest
@AutoConfigureTestDatabase(replace = AutoConfigureTestDatabase.Replace.NONE)
class ItemControllerTest {

    @Autowired
    private TestRestTemplate restTemplate;

    @Autowired
    private ItemRepository itemRepository;

    @Test
    void shouldCreateItem() {
        // Given
        Item newItem = new Item();
        newItem.setName("Test Item");
        newItem.setDescription("Test Description");
        newItem.setCategory("Electronics");

        // When
        ResponseEntity<Item> response = restTemplate.postForEntity(
            "/api/items", newItem, Item.class);

        // Then
        assertThat(response.getStatusCode()).isEqualTo(HttpStatus.CREATED);
        assertThat(response.getBody().getName()).isEqualTo("Test Item");

        // Verify in database
        Optional<Item> savedItem = itemRepository.findById(response.getBody().getId());
        assertThat(savedItem).isPresent();
    }

    @Test
    void shouldReturnValidationErrorForInvalidItem() {
        // Given
        Item invalidItem = new Item();
        // Name ist leer -> Validierungsfehler

        // When
        ResponseEntity<String> response = restTemplate.postForEntity(
            "/api/items", invalidItem, String.class);

        // Then
        assertThat(response.getStatusCode()).isEqualTo(HttpStatus.BAD_REQUEST);
    }
}
\end{lstlisting}

\section{Deployment}

\subsection{Docker}

\textbf{Warum Docker?}
\begin{itemize}
    \item \textbf{Konsistenz:} "It works on my machine" Problem lösen
    \item \textbf{Isolation:} Anwendungen laufen in eigenen Containern
    \item \textbf{Skalierung:} Container können einfach repliziert werden
    \item \textbf{Portabilität:} Läuft überall wo Docker läuft
\end{itemize}

\begin{lstlisting}[language=Docker, caption=Dockerfile für Spring Boot Backend]
FROM openjdk:17-jdk-slim

WORKDIR /app

# Maven Dependencies kopieren (für besseres Caching)
COPY pom.xml .
COPY mvnw .
COPY .mvn .mvn
RUN ./mvnw dependency:go-offline

# Source Code kopieren
COPY src ./src

# Application bauen
RUN ./mvnw clean package -DskipTests

# JAR-File ausführen
EXPOSE 8080
CMD ["java", "-jar", "target/backend-0.0.1-SNAPSHOT.jar"]
\end{lstlisting}

\subsection{Docker Compose}

\begin{lstlisting}[language=YAML, caption=docker-compose.yml]
version: '3.8'

services:
  database:
    image: postgres:15
    environment:
      POSTGRES_DB: cictogive
      POSTGRES_USER: cicuser
      POSTGRES_PASSWORD: cicpass
    ports:
      - "5432:5432"
    volumes:
      - postgres_data:/var/lib/postgresql/data

  backend:
    build: ./backend
    ports:
      - "8080:8080"
    environment:
      - SPRING_DATASOURCE_URL=jdbc:postgresql://database:5432/cictogive
      - SPRING_DATASOURCE_USERNAME=cicuser
      - SPRING_DATASOURCE_PASSWORD=cicpass
    depends_on:
      - database

  frontend:
    build: ./frontend
    ports:
      - "3000:3000"
    environment:
      - NEXT_PUBLIC_API_URL=http://localhost:8080
    depends_on:
      - backend

volumes:
  postgres_data:
\end{lstlisting}

\section{Performance-Optimierung}

\subsection{Backend-Performance}

\begin{itemize}
    \item \textbf{Datenbankindizes:} Auf häufig gesuchte Felder
    \item \textbf{Connection Pooling:} Datenbankverbindungen wiederverwenden
    \item \textbf{Caching:} Redis für häufig abgerufene Daten
    \item \textbf{Pagination:} Große Datensätze aufteilen
    \item \textbf{Lazy Loading:} Related Entities nur bei Bedarf laden
\end{itemize}

\subsection{Frontend-Performance}

\begin{itemize}
    \item \textbf{Code Splitting:} Nur benötigter Code wird geladen
    \item \textbf{Image Optimization:} Next.js Image-Komponente
    \item \textbf{Caching:} Browser-Cache und CDN
    \item \textbf{Bundle Analysis:} Große Dependencies identifizieren
    \item \textbf{Lazy Loading:} Komponenten und Bilder bei Bedarf laden
\end{itemize}

\section{Projektmanagement und Zeitschätzung}

\subsection{Agile Entwicklung}

\textbf{Scrum/Kanban Prinzipien:}
\begin{itemize}
    \item \textbf{Iterative Entwicklung:} In kleinen Schritten entwickeln
    \item \textbf{MVP (Minimum Viable Product):} Erst Grundfunktionen, dann erweitern
    \item \textbf{User Stories:} Anforderungen aus Benutzersicht
    \item \textbf{Daily Standups:} Regelmäßiger Austausch (bei Teams)
\end{itemize}

\subsection{Detaillierte Zeitschätzung}

\begin{center}
\begin{tabularx}{\textwidth}{|X|c|X|}
\hline
\textbf{Task} & \textbf{Stunden} & \textbf{Begründung} \\
\hline
Projektsetup (Maven, Node.js, Git) & 2 & Grundkonfiguration \\
\hline
Datenmodell Design \& Implementierung & 8 & 3 Entities + Relationen \\
\hline
Spring Security + JWT & 12 & Komplex, viele Konfigurationen \\
\hline
User Registration/Login & 6 & Backend + Frontend \\
\hline
Item CRUD Operations & 10 & Alle HTTP-Methoden \\
\hline
File Upload System & 8 & Validierung + Storage \\
\hline
Search \& Filter Logic & 8 & Database Queries + Frontend \\
\hline
Comment System & 6 & Relations + UI \\
\hline
Frontend Components & 20 & Alle React-Komponenten \\
\hline
API Integration & 8 & Fetch Calls + Error Handling \\
\hline
Responsive Design & 12 & Tailwind CSS + Mobile \\
\hline
Testing & 16 & Unit + Integration Tests \\
\hline
Deployment Setup & 6 & Docker + Environment \\
\hline
Bug Fixing \& Polish & 12 & Unvorhergesehene Probleme \\
\hline
Documentation & 8 & Code Comments + README \\
\hline
\textbf{Total} & \textbf{142} & \textbf{≈ 18 Arbeitstage} \\
\hline
\end{tabularx}
\end{center}

\subsection{Risikomanagement}

\textbf{Häufige Probleme und Lösungen:}

\begin{itemize}
    \item \textbf{CORS-Probleme:} Früh testen, richtige Headers setzen
    \item \textbf{JWT-Token-Handling:} Expiration-Zeiten beachten
    \item \textbf{File Upload-Limits:} Server und Frontend-Limits abstimmen
    \item \textbf{Database Schema Changes:} Migration-Strategien planen
    \item \textbf{Mobile Responsiveness:} Früh auf verschiedenen Geräten testen
\end{itemize}

\section{Fazit und Ausblick}

\subsection{Was wurde erreicht?}

Das Projekt "CIC to Give" demonstriert eine vollständige moderne Webentwicklung:

\textbf{Backend-Kompetenzen:}
\begin{itemize}
    \item Spring Boot Framework-Verständnis
    \item RESTful API-Design
    \item Datenbank-Modellierung und JPA
    \item Sicherheitskonzepte (JWT, BCrypt)
    \item File Upload und Validation
\end{itemize}

\textbf{Frontend-Kompetenzen:}
\begin{itemize}
    \item React/Next.js Entwicklung
    \item TypeScript für Typsicherheit
    \item Responsive Design mit Tailwind CSS
    \item API-Integration und State Management
    \item User Experience Design
\end{itemize}

\textbf{DevOps-Kompetenzen:}
\begin{itemize}
    \item Docker Container-Technologie
    \item Database Setup und Configuration
    \item Testing-Strategien
    \item Deployment-Prozesse
\end{itemize}

\subsection{M{\"o}gliche Erweiterungen}

\begin{itemize}
    \item \textbf{Real-time Chat:} WebSocket f{\"u}r Kommunikation zwischen Nutzern
    \item \textbf{Push Notifications:} Benachrichtigungen f{\"u}r neue Items
    \item \textbf{Geolocation:} Karten-Integration für Standorte
    \item \textbf{Rating System:} Bewertungen f{\"u}r Nutzer
    \item \textbf{Admin Panel:} Content-Moderation
    \item \textbf{Mobile App:} React Native oder Flutter
\end{itemize}

\subsection{Lessons Learned}

\textbf{Technische Erkenntnisse:}
\begin{itemize}
    \item Fr{\"u}he Integration zwischen Frontend und Backend spart Zeit
    \item Gute Dokumentation und Typisierung erleichtern Entwicklung
    \item Testing ist Investment in zuk{\"u}nftige Wartbarkeit
    \item Security sollte von Anfang an mitgedacht werden
\end{itemize}

\textbf{Projektmanagement:}
\begin{itemize}
    \item Realistische Zeitsch{\"a}tzungen sind schwierig aber wichtig
    \item MVP-Ansatz hilft bei der Priorisierung
    \item Regelm{\"a}{\ss}ige Deployments decken Probleme fr{\"u}h auf
\end{itemize}

\vfill
\begin{center}
\textit{Diese Dokumentation zeigt die vollst{\"a}ndige Entwicklung einer modernen Fullstack-Anwendung von der Konzeption bis zum Deployment.}

\vspace{1cm}

\textbf{Lehrabschlusspr{\"u}fung Applikationsentwicklung - Coding}\\
\textit{Dokument erstellt am: \today}\\
\textit{Version: 1.0}
\end{center}

\end{document}
}
